\documentclass[12pt,a4paper]{article}
\usepackage[spanish]{babel}
\usepackage[utf8]{inputenc}
\usepackage{array}
\usepackage{longtable}
\usepackage{tcolorbox}
\usepackage{hyperref}
\usepackage{geometry} % Paquete para márgenes

% Márgenes más amplios (ajústalos a tu gusto)
\geometry{
  top=3cm,
  bottom=3cm,
  left=3cm,
  right=3cm
}

\tcbset{
  mydialog/.style={
    colback=black!90,
    coltext=yellow!90,
    colframe=yellow,
    boxrule=1pt,
    arc=8pt,
    left=10pt, right=10pt, top=10pt, bottom=10pt,
    fonttitle=\bfseries,
    coltitle=yellow,
    title=Información de contacto,
    enhanced,
    drop shadow,
    borderline={0.5pt}{0pt}{yellow}
  }
}

\title{TRABAJO FINAL:\\ PROYECTO DE TESIS\\ Seminario II}
\author{Mtro. Staling Cordero-Brito\\
Historia y Geografía para Educadores}

\date{\today}

\begin{document}
\maketitle

\section*{Rúbrica de Evaluación (40 puntos)}

\renewcommand{\arraystretch}{1.5}
\begin{longtable}{|>{\raggedright}p{9cm}|c|}
\hline
\textbf{Criterios de Evaluación} & \textbf{Puntos} \\
\hline
\multicolumn{2}{|l|}{\textbf{Nota previa}} \\
Los \textbf{Resumen (palabras clave) e Introducción (con asesor)} deberán incluirse siguiendo las Guías suministradas, 
pero quedarán \textbf{como borrador}. & -- \\
\hline
\multicolumn{2}{|l|}{\textbf{1. Planteamiento del problema (15 pts)}} \\
Claridad en formulación del problema, delimitación, justificación, objetivos y preguntas de investigación. & 15 \\
\hline
\multicolumn{2}{|l|}{\textbf{2. Marco Teórico (12 pts)}} \\
\textbf{2.1 Relevancia y coherencia con el problema de investigación} & 2 \\
\textbf{2.2 Calidad y actualidad de las fuentes utilizadas} & 2 \\
\textbf{2.3 Profundidad en la construcción de categorías analíticas} & 2 \\
\textbf{2.4 Capacidad crítica y comparativa } & 2 \\
\textbf{2.5 Articulación con el enfoque} & 2 \\
\textbf{2.6 Pertinencia con el tema de investigación} & 2 \\
\hline
\multicolumn{2}{|l|}{\textbf{3. Metodología (5 pts)}} \\
Coherencia entre tipo de investigación; enfoque; diseño; alcance; métodos; técnicas; instrumentos; universo, población y muestra. & 5 \\
\hline
\multicolumn{2}{|l|}{\textbf{4. Presentación de Resultados (3 pts)}} \\
Se presentarán los resultados esperados para poner en contexto al asesor. & 3 \\
\hline
\multicolumn{2}{|l|}{\textbf{5. Discusión y Conclusiones (con asesor)}} \\
Este apartado quedará en borrador, a ser desarrollado con el asesor. & -- \\
\hline
\multicolumn{2}{|l|}{\textbf{6. Bibliografía (5 pts)}} \\
Debe presentarse en estilo Chicago, siguiendo la Guía suministrada. & 5 \\
\hline
\multicolumn{1}{|r|}{\textbf{Total}} & \textbf{40} \\
\hline
\end{longtable}


\begin{tcolorbox}[
    enhanced,
    frame hidden,
    width=0.8\textwidth,
    boxrule=0pt,
    colback=gray!20,
    sharp corners,
    borderline={0.5pt}{0pt}{black!70},
    drop shadow
  ]

  % Parte superior (gris oscuro)
  \begin{tcolorbox}[
      colback=gray!40,
      coltext=white,
      boxrule=0pt,
      sharp corners,
      left=5pt,right=5pt,top=5pt,bottom=5pt
    ]
    \centering
    \textbf{Mtro. Staling Cordero-Brito}
  \end{tcolorbox}

  % Parte inferior (gris claro)
  \begin{tcolorbox}[
      colback=gray!15,
      coltext=black,
      boxrule=0pt,
      sharp corners,
      left=5pt,right=5pt,top=5pt,bottom=5pt
    ]
    \centering
    Maestría en Historia y Geografía para Educadores \\[4pt]
    Universidad Autónoma de Santo Domingo (UASD) \\[4pt]
    Correo: \texttt{scordero49@uasd.edu.do}
  \end{tcolorbox}

\end{tcolorbox}

\end{document}
